\chapter{Subsystem integration}
\label{ch:integration}

\section{Shared AXI4-Lite slave}
\label{sec:axislave}

\file{hdl/axi\_lite\_slave.v} implements the slave-side handshake once. Both the
PIC and the machine timer instantiate it, and it is available unchanged to any
register-mapped peripheral --- the DP-SRAM ports or the DMA configuration port
would use it as-is.

The split is deliberate: the shared module owns AW/W collection, the B response
and the AR/R response, while the peripheral only describes its registers. A
peripheral supplies four things:

\begin{table}[H]
\centering
\caption{Peripheral-side hooks of \file{axi\_lite\_slave.v}}
\label{tab:slavehooks}
\begin{tabular}{@{}L{0.22\linewidth}lL{0.46\linewidth}@{}}
\toprule
\textbf{Hook} & \textbf{Direction} & \textbf{Meaning} \\
\midrule
\sig{wr\_en}, \sig{wr\_addr}, \sig{wr\_data}, \sig{wr\_strb}
  & to peripheral & One-cycle write pulse with word offset, data and strobes. \\
\sig{wr\_ok}   & from peripheral & Combinational: is this offset writable? \\
\sig{rd\_addr} & to peripheral & Requested word offset, valid while \sig{ARVALID}. \\
\sig{rd\_data}, \sig{rd\_ok} & from peripheral & Combinational read multiplexer and readability qualifier. \\
\bottomrule
\end{tabular}
\end{table}

\subsection{Write path}

AW and W are collected independently and may arrive in either order. Two flags
track what has landed; when both are present and no response is outstanding, the
transaction commits.

\begin{center}
\ttfamily\small
wr\_commit = aw\_got\_q \&\& w\_got\_q \&\& !s\_axi\_bvalid\_o;
\end{center}

Readiness is derived directly from those flags, so the slave stops accepting a
channel it has already captured and stops accepting anything at all while a
response is pending:

\begin{center}
\ttfamily\small
s\_axi\_awready\_o = !aw\_got\_q \&\& !s\_axi\_bvalid\_o;\\
s\_axi\_wready\_o\ \ = !w\_got\_q\ \ \&\& !s\_axi\_bvalid\_o;
\end{center}

\begin{figure}[H]
\centering
\begin{tikzpicture}[x=1mm,y=1mm,>={Stealth[length=2.4mm]},semithick]
  \node[stbi] (idle)  at (0,0)    {IDLE};
  \node[stb]  (gotaw) at (46,20)  {GOT\_AW};
  \node[stb]  (gotw)  at (46,-20) {GOT\_W};
  \node[stb]  (both)  at (92,0)   {BOTH};
  \node[stb]  (resp)  at (137,0)  {RESP};

  \draw[->,draw=accent] (idle)  -- node[lbl,pos=0.5,above left=-1mm and -3mm]{AW hs} (gotaw);
  \draw[->,draw=accent] (idle)  -- node[lbl,pos=0.5,below left=-1mm and -3mm]{W hs}  (gotw);
  \draw[->,draw=accent] (gotaw) -- node[lbl,pos=0.5,above right=-1mm and -3mm]{W hs} (both);
  \draw[->,draw=accent] (gotw)  -- node[lbl,pos=0.5,below right=-1mm and -3mm]{AW hs}(both);
  \draw[->,draw=accent] (idle)  -- node[lbl,pos=0.5,above,align=center]
        {AW + W\\\scriptsize same cycle} (both);
  \draw[->,draw=accent] (both)  -- node[lbl,pos=0.5,above,align=center]
        {\texttt{wr\_commit}\\\scriptsize \texttt{wr\_en\_o} pulses} (resp);

  % return leg routed below every node, clear of GOT_W
  \draw[->,draw=accent] (resp.south) -- (137,-38) -- (0,-38) -- (idle.south);
  \node[lbl,above=0.5mm] at (68,-38) {BREADY --- BVALID drops, flags clear};

  \node[align=left,font=\scriptsize,anchor=north west,
        draw=rulegrey,fill=boxbg,inner sep=5pt,text width=118mm] at (0,-48)
    {\textbf{Reset:} both flags clear, \texttt{BVALID} low.
     \texttt{BOTH} lasts exactly one cycle --- the got bits are registered and
     \texttt{wr\_commit} is combinational from them, so the commit is always the
     cycle \emph{after} the second handshake, including when AW and W arrive
     together. \texttt{BRESP} = \texttt{wr\_ok\_i}\,?\,OKAY\,:\,SLVERR, captured
     at commit.};
\end{tikzpicture}
\caption{Write handshake of \file{axi\_lite\_slave.v}, encoded in
\sig{aw\_got\_q}, \sig{w\_got\_q} and \sig{s\_axi\_bvalid\_o} rather than in a
state register.}
\label{fig:slavewr}
\end{figure}

\subsection{Read path}

The read side is simpler: \sig{ARREADY} is the inverse of \sig{RVALID}, so a new
address is accepted only when no response is pending. On the AR handshake the
peripheral's combinational read data and its \sig{rd\_ok} qualifier are captured
into \sig{RDATA} and \sig{RRESP}.

\sig{BRESP} and \sig{RRESP} are combinational functions of stored error flags
and are driven whether or not the corresponding \sig{VALID} is asserted, so both
flags and \sig{s\_axi\_rdata\_o} reset to zero: the response pins read
\texttt{OKAY} and \texttt{0x0000\_0000} from time zero rather than \texttt{X}.
The reset bench checks all three while reset is still asserted
(Section~\ref{sec:vp-picreset}, step 1).

\begin{figure}[H]
\centering
\begin{tikztimingtable}[timing/slope=0,timing/coldist=2pt,xscale=1.05]
  clk      & 10{C} \\
  ARVALID  & 2L 2H 6L \\
  ARREADY  & 4H 2L 4H \\
  ARADDR   & 2U 2D{off} 6U \\
  RVALID   & 4L 4H 2L \\
  RREADY   & 6L 2H 2L \\
  RDATA    & 4U 4D{reg} 2U \\
  \extracode
  \begin{pgfonlayer}{background}
    \vertlines[help lines]{2,4,6,8}
  \end{pgfonlayer}
\end{tikztimingtable}
\caption{Register read. \sig{ARREADY} drops while the response is outstanding,
which is what enforces the one-transaction-per-direction contract.}
\label{fig:slaveread}
\end{figure}

\paragraph{Contract.} At most one transaction per direction, which matches the
CPU's one-outstanding master exactly. Nothing in the system requires more.

\section{Interconnect in the verification environment}

\begin{sloppypar}
The testbench builds the interconnect from two cascaded 1-to-2 address decoders,
\file{axi\_lite\_dec2.v}. The first splits the \texttt{0x3xxx\_xxxx} window from
everything else; the second splits that window into the PIC and the machine
timer. Anything outside a mapped region returns DECERR from the memory model,
which the CPU turns into a precise access fault.
\end{sloppypar}

This is testbench infrastructure, not part of the delivery. The real system will
supply its own fabric; the CPU requires nothing from it beyond conformant
AXI4-Lite behaviour and the three response codes.

\section{Interrupt path end to end}

\begin{figure}[H]
\centering
\begin{tikzpicture}[node distance=4mm,font=\small]
  \node[blkf,minimum width=26mm] (p) {peripheral};
  \node[blk,right=14mm of p,minimum width=26mm] (pic) {pic};
  \node[blk,right=14mm of pic,minimum width=26mm] (cpu) {cpu\_top};
  \node[blkf,right=14mm of cpu,minimum width=26mm] (sw) {handler};

  \draw[flow] (p) -- node[lbl,above,align=center]{\scriptsize line asserted\\\scriptsize and held} (pic);
  \draw[flow] (pic) -- node[lbl,above,align=center]{\scriptsize cpu\_irq\\\scriptsize + vector} (cpu);
  \draw[flow] (cpu) -- node[lbl,above,align=center]{\scriptsize trap entry\\\scriptsize mcause = 16+v} (sw);
  \draw[flow] (cpu.south) ++(-6mm,0) -- ++(0,-8mm) -| node[lbl,pos=0.25,below]{\scriptsize ack: claim, push} (pic.south);
  \draw[flow] (sw.south) -- ++(0,-14mm) -| node[lbl,pos=0.25,below]{\scriptsize clears the peripheral} (p.south);
  \draw[flow] (sw.north) -- ++(0,6mm) -| node[lbl,pos=0.25,above]{\scriptsize MRET $\to$ eoi: pop} (pic.north);
\end{tikzpicture}
\caption{Interrupt flow across the three blocks. The peripheral holds its line
until the handler clears it; the PIC owns priority and nesting; the CPU owns the
architectural trap state.}
\label{fig:irqflow}
\end{figure}

The two masks stack on purpose. \reg{INT\_ENABLE} in the PIC says ``this source
may reach the CPU''; \reg{mie} in the CPU says ``current software cares''. They
are independent, which is what allows software to mask a class of interrupts
without reconfiguring the controller.

Stacking them naively does not work, and the first version of this design got it
wrong. The PIC resolved the winner without knowing \reg{mie}, and the core then
applied \reg{mie} to the winner it had been handed. A source enabled in the PIC
but masked in \reg{mie} was therefore selected and never claimed: the controller
parked on it and every less urgent source starved behind it, and a \sig{WFI}
waiting on one of those sources never woke. It was written up at the time as an
integration rule --- do not enable a source no handler will service --- but a rule
that forbids using a per-source mask for what it is for is not a rule, it is a
defect with a workaround.

The core now exports \reg{mie}[31:16] to the PIC as \sig{cpu\_mask\_i}, and the
resolver skips a masked source when it chooses. The source stays pending, so its
deadline counter keeps running and it is offered the moment software unmasks it;
the core still re-checks the mask on the offer it receives, which covers the
cycle in which \reg{mie} changes. Both masks now do what their names say.

\section{Multi-core behaviour}

Nothing in the core is shared between instances. There is no static state,
memories are external, and identity comes only from the \sig{RESET\_PC} and
\sig{HART\_ID} parameters. Two or more instances work behind a standard AXI
interconnect.

RV32I has no load-reserved / store-conditional, but every memory operation here
is in-order and blocking --- one outstanding transaction, no store buffer --- so
each hart is sequentially consistent and flag-based handshakes through shared
memory are sound. This is the same property the system's ping-pong buffering
relies on.

\begin{deviation}
\textbf{Not resolved at system level.} The PIC serves one CPU. Which source
targets which core is undefined for a multi-core configuration, and each core
would want its own timer compare register. The dual-core testbench exercises two
cores sharing memory through an arbiter; it does not exercise interrupt routing
to two cores, because the routing policy does not exist yet.
\end{deviation}
